%% ma: L12-B08-0047
%% lop: 12
%% chuong: 2
%% bai: 8
%% dang: 12.8.12
%% loai: TN4
%% muc: TH
%% dap_an: "A"
%% nguon: "LaTeX giáo viên gửi 21/10/2026 (M05, Vectơ và hệ trục tọa độ trong không gian, 107 câu), Phần I Câu 54"
%% kien_thuc: "Bài 6 (quy tắc hình bình hành), Bài 8 (hai vectơ bằng nhau theo tọa độ)"
%% trang_thai: chinh-thuc
%% ghi_chu: "Trùng ý tưởng với L12-B08-0017 (tìm đỉnh hình bình hành), mở rộng cho hình hộp, nên nhập cùng dạng 12.8.12. Phần I Câu 20 và Câu 37 (tìm đỉnh D của hình bình hành) chỉ đổi số so với L12-B08-0017 nên đã bỏ"
\begin{ex}
Trong không gian $Oxyz$, cho hình hộp $ABCD.A'B'C'D'$ với $A(1;2;0)$, $B(3;2;0)$, $D(1;5;0)$, $C'(4;6;-2)$. Tọa độ đỉnh $A'$ là
\choice{\True $(2;3;-2)$}{$(2;3;2)$}{$(4;6;-4)$}{$(4;6;-2)$}
\loigiai{$ABCD$ là hình bình hành nên $\overrightarrow{AC}=\overrightarrow{AB}+\overrightarrow{AD}=(2;3;0)$, suy ra $C(3;5;0)$.

$\overrightarrow{AA'}=\overrightarrow{CC'}=(1;1;-2)$ nên $A'(2;3;-2)$.}
\end{ex}
